\section{Probability Model}\label{probability-model}

We derive the relationship between DOF and error probability from \texttt{AbstractClassSystem}'s axis independence theorem.

\subsection{Error Independence from Axis Orthogonality}

The independence of errors is not an axiom; it is a consequence of axis orthogonality proved in \texttt{AbstractClassSystem}~\cite{paper1_typing_discipline}.

\begin{theorem}[Error Independence]\label{thm:error-independence}
If axes $\{A_1, \ldots, A_n\}$ are orthogonal (\texttt{AbstractClassSystem}, Theorem \texttt{minimal\_complete\_unique\_orthogonal}), then errors along each axis are statistically independent.
\end{theorem}
\leanmeta{\LH{L23}}

\begin{proof}
By Definition~\ref{def:architecture}, the state space is the product $S = \prod_{c \in C} S_c$. \texttt{AbstractClassSystem}'s orthogonality theorem (\texttt{minimal\_complete\_unique\_orthogonal}) establishes that a minimal complete axis set partitions the state space into independent dimensions: axis $A_i$ is a projection onto a factor of the product, and $A_i \perp A_j$ means the projection onto factor $i$ carries no information about factor $j$.

An \emph{error along axis $A_i$} is the event $E_i$ that the value on factor $i$ deviates from specification. Since the factors are independent dimensions of the product space, the probability measure factorizes:
\[
P(E_i \mid \text{state of factors } \neq i) = P(E_i).
\]
This is the standard characterization of independence for product probability spaces: the marginal on factor $i$ is independent of the marginal on any other factor. Therefore $P(E_i \cap E_j) = P(E_i) \cdot P(E_j)$ for all $i \neq j$, and by induction the $n$ events $E_1, \ldots, E_n$ are mutually independent.
\end{proof}

\begin{corollary}[DOF = Independent Error Sources]\label{cor:dof-errors}
DOF$(A) = n$ implies $n$ independent sources of error, each with probability $p$.
\end{corollary}
\leanmeta{\LH{L15}, \LH{L23}}

\begin{proof}
DOF counts independent axes (Section~\ref{foundations}, Definition~\ref{def:dof}). By Theorem~\ref{thm:error-independence}, independent axes have independent errors.
\end{proof}

\begin{theorem}[Error Compounding]\label{thm:error-compound}
For a system to be correct, all $n$ independent axes must be error-free. Errors compound multiplicatively.
\end{theorem}
\leanmeta{\LH{L20}, \LH{L41}}

\begin{proof}
By \texttt{Ssot}'s coherence theorem (\texttt{oracle\_arbitrary}), incoherence in any axis violates system correctness. An error in axis $A_i$ introduces incoherence along $A_i$. Therefore, correctness requires $\bigwedge_{i=1}^{n} \neg\text{error}(A_i)$. By Theorem~\ref{thm:error-independence}, this probability is $(1-p)^n$.
\end{proof}

\subsection{Error Probability Formula}

\begin{theorem}[Error Probability]\label{thm:error-prob}
For architecture with $n$ DOF and per-component error rate $p$:
\[
P_{\text{error}}(n) = 1 - (1-p)^n
\]
\end{theorem}
\leanmeta{\LH{L24}, \LH{L40}}

\begin{proof}
By Theorem~\ref{thm:error-independence} (derived from \texttt{AbstractClassSystem}'s orthogonality), each DOF has independent error probability $p$, so each is correct with probability $(1-p)$. By Theorem~\ref{thm:error-compound}, all $n$ DOF must be correct:
\[
P_{\text{correct}}(n) = (1-p)^n
\]
Therefore:
\[
P_{\text{error}}(n) = 1 - P_{\text{correct}}(n) = 1 - (1-p)^n
\]
\end{proof}

\begin{corollary}[Linear Approximation]\label{cor:linear-approx}
For fixed $p>0$, the linear expected-error model and the exact series model induce the same ordering on architectures by DOF.
\end{corollary}
\leanmeta{\LH{L16}, \LH{L39}}

\begin{proof}
By Theorem~\ref{thm:approx-bound}, ordering equivalence is exact under the discrete model used in the mechanization.
\end{proof}

\begin{corollary}[DOF-Error Monotonicity]\label{cor:dof-monotone}
For architectures $A_1, A_2$:
\[
\text{DOF}(A_1) < \text{DOF}(A_2) \implies P_{\text{error}}(A_1) < P_{\text{error}}(A_2)
\]
\end{corollary}
\leanmeta{\LH{L33}}

\begin{proof}
$P_{\text{error}}(n) = 1 - (1-p)^n$ is strictly increasing in $n$ for $p \in (0,1)$.
\end{proof}

\subsection{Expected Errors}

\begin{theorem}[Expected Error Bound]\label{thm:expected-errors}
Expected number of errors in architecture $A$:
\[
\mathbb{E}[\text{\# errors}] = p \cdot \text{DOF}(A)
\]
\end{theorem}
\leanmeta{\LHrng{L}{25}{26}}

\begin{proof}
By linearity of expectation:
\[
\mathbb{E}[\text{\# errors}] = \sum_{i=1}^{\text{DOF}(A)} P(\text{error in DOF}_i) = \sum_{i=1}^{\text{DOF}(A)} p = p \cdot \text{DOF}(A)
\]
\end{proof}

\begin{example}[Concrete Calculations]
Assume $p = 0.01$ (1\% per-component error rate):
\begin{itemize}
\item DOF $= 1$: $P_{\text{error}} = 1 - 0.99 = 0.01$ (1\%)
\item DOF $= 10$: $P_{\text{error}} = 1 - 0.99^{10} \approx 0.096$ (9.6\%)
\item DOF $= 100$: $P_{\text{error}} = 1 - 0.99^{100} \approx 0.634$ (63.4\%)
\end{itemize}
\end{example}

\subsection{Connection to Reliability Theory}

The error model has a direct interpretation in classical reliability theory \cite{patterson2013computer}, connecting software architecture to a mature mathematical framework with 60+ years of theoretical development.

\begin{theorem}[DOF-Reliability Isomorphism]\label{thm:dof-reliability}
An architecture with DOF $= n$ is \emph{isomorphic} to a series reliability system with $n$ components. The isomorphism:
\begin{enumerate}
\item \textbf{Preserves ordering:} If $\text{DOF}(A_1) < \text{DOF}(A_2)$, then $P_{\text{error}}(A_1) < P_{\text{error}}(A_2)$
\item \textbf{Is invertible:} Round-trip mapping preserves DOF exactly
\item \textbf{Connects domains:} $P_{\text{error}}(n) = 1 - R_{\text{series}}(n)$ where $R_{\text{series}}(n) = (1-p)^n$
\end{enumerate}
\end{theorem}
\leanmeta{\LH{L22}, \LH{L27}}

\textbf{Interpretation:} Each DOF is a ``component'' that must work correctly. This is the reliability analog of Theorem~\ref{thm:error-independence}, which derives error independence from axis orthogonality.

\begin{theorem}[Ordering Equivalence (Exact)]\label{thm:approx-bound}
For architectures $A_1,A_2$ with equal capabilities and error rate $p>0$, the linear model and the exact series model induce the same DOF ordering:
\[
\text{DOF}(A_1)\le \text{DOF}(A_2)
\iff
(\mathbb{E}[\text{errors}(A_1)]\cdot d)\le(\mathbb{E}[\text{errors}(A_2)]\cdot d),
\]
where $d$ is the denominator of the discrete error-rate representation.
\end{theorem}
\leanmeta{\LH{L32}, \LH{L39}}

\begin{proof}
This is theorem \texttt{ordering\_equivalence\_exact} in \texttt{Leverage/Probability.lean}.
\end{proof}

\textbf{Linear Approximation Justification:} For small $p$ (the software engineering regime where $p \approx 0.01$), the linear model $P_{\text{error}} \approx n \cdot p$ is:
\begin{enumerate}
\item Order-equivalent to the exact series model in the mechanized discrete representation
\item Monotone in DOF for fixed positive error rate
\item Cleanly formalized in natural-number arithmetic
\end{enumerate}

\subsection{Epistemic Grounding}

The probability model is not axiomatic; it is derived from the epistemic foundations in \texttt{AbstractClassSystem} and \texttt{Ssot}:

\begin{enumerate}
\item \texttt{AbstractClassSystem} proves axis orthogonality (\texttt{minimal\_complete\_unique\_orthogonal})
\item \textbf{Theorem~\ref{thm:error-independence}} derives error independence from orthogonality
\item \texttt{Ssot} establishes that DOF = 1 guarantees coherence (\texttt{oracle\_arbitrary})
\item \textbf{Theorem~\ref{thm:error-compound}} connects errors to incoherence
\end{enumerate}

This derivation chain ensures the probability model rests on proved foundations, not assumed axioms.

\subsection{Leverage Gap: Quantitative Predictions}

The leverage framework provides \emph{quantitative}, empirically testable predictions about modification costs.

\begin{theorem}[Leverage Gap]\label{thm:leverage-gap}
For architectures with equal capabilities, the modification cost ratio equals the inverse leverage ratio:
\[
\frac{\text{ModCost}(A_2)}{\text{ModCost}(A_1)} = \frac{\text{DOF}(A_2)}{\text{DOF}(A_1)} = \frac{L(A_1)}{L(A_2)}
\]
\end{theorem}
\leanmeta{\LH{L30}}

\begin{theorem}[Testable Prediction]\label{thm:testable-prediction}
If architecture $A_1$ has $n\times$ lower DOF than $A_2$ (for equal capabilities), then $A_1$ requires $n\times$ fewer expected modifications. This is empirically testable against PR/commit data.
\end{theorem}
\leanmeta{\LH{L42}}

\begin{corollary}[DOF Ratio Predicts Error Ratio]\label{cor:dof-ratio}
The ratio of expected errors between two architectures equals the ratio of their DOF:
\[
\frac{\mathbb{E}[\text{errors}(A_2)]}{\mathbb{E}[\text{errors}(A_1)]} = \frac{\text{DOF}(A_2)}{\text{DOF}(A_1)}
\]
\end{corollary}
\leanmeta{\LH{L21}}

These theorems transform architectural intuitions into testable hypotheses. A claim that ``Pattern X is 3× better than Pattern Y'' can be verified by comparing DOF and measuring modification frequency in real codebases.

\subsection{Formalization}

Formalized in \texttt{Leverage/Probability.lean}:
\begin{itemize}
\item \texttt{dof\_reliability\_isomorphism}: Theorem~\ref{thm:dof-reliability} (the central isomorphism)
\item \texttt{isomorphism\_preserves\_failure\_ordering}: Ordering preservation
\item \texttt{isomorphism\_roundtrip}: Invertibility proof
\item \texttt{ordering\_equivalence\_exact}: Theorem~\ref{thm:approx-bound} (exact ordering equivalence)
\item \texttt{linear\_model\_preserves\_ordering}: Linear ordering support
\item \texttt{leverage\_gap}: Theorem~\ref{thm:leverage-gap}
\item \texttt{testable\_modification\_prediction}: Theorem~\ref{thm:testable-prediction}
\item \texttt{dof\_ratio\_predicts\_error\_ratio}: Corollary~\ref{cor:dof-ratio}
\item \texttt{lower\_dof\_lower\_errors}: Corollary~\ref{cor:dof-monotone}
\item \texttt{expected\_errors\_from\_linearity}: Theorem~\ref{thm:expected-errors}
\item \texttt{ssot\_minimal\_errors}: SSOT minimality
\end{itemize}
